% Options for packages loaded elsewhere
\PassOptionsToPackage{unicode}{hyperref}
\PassOptionsToPackage{hyphens}{url}
\documentclass[
]{article}
\usepackage{xcolor}
\usepackage[margin=1in]{geometry}
\usepackage{amsmath,amssymb}
\setcounter{secnumdepth}{-\maxdimen} % remove section numbering
\usepackage{iftex}
\ifPDFTeX
  \usepackage[T1]{fontenc}
  \usepackage[utf8]{inputenc}
  \usepackage{textcomp} % provide euro and other symbols
\else % if luatex or xetex
  \usepackage{unicode-math} % this also loads fontspec
  \defaultfontfeatures{Scale=MatchLowercase}
  \defaultfontfeatures[\rmfamily]{Ligatures=TeX,Scale=1}
\fi
\usepackage{lmodern}
\ifPDFTeX\else
  % xetex/luatex font selection
\fi
% Use upquote if available, for straight quotes in verbatim environments
\IfFileExists{upquote.sty}{\usepackage{upquote}}{}
\IfFileExists{microtype.sty}{% use microtype if available
  \usepackage[]{microtype}
  \UseMicrotypeSet[protrusion]{basicmath} % disable protrusion for tt fonts
}{}
\makeatletter
\@ifundefined{KOMAClassName}{% if non-KOMA class
  \IfFileExists{parskip.sty}{%
    \usepackage{parskip}
  }{% else
    \setlength{\parindent}{0pt}
    \setlength{\parskip}{6pt plus 2pt minus 1pt}}
}{% if KOMA class
  \KOMAoptions{parskip=half}}
\makeatother
\usepackage{color}
\usepackage{fancyvrb}
\newcommand{\VerbBar}{|}
\newcommand{\VERB}{\Verb[commandchars=\\\{\}]}
\DefineVerbatimEnvironment{Highlighting}{Verbatim}{commandchars=\\\{\}}
% Add ',fontsize=\small' for more characters per line
\usepackage{framed}
\definecolor{shadecolor}{RGB}{248,248,248}
\newenvironment{Shaded}{\begin{snugshade}}{\end{snugshade}}
\newcommand{\AlertTok}[1]{\textcolor[rgb]{0.94,0.16,0.16}{#1}}
\newcommand{\AnnotationTok}[1]{\textcolor[rgb]{0.56,0.35,0.01}{\textbf{\textit{#1}}}}
\newcommand{\AttributeTok}[1]{\textcolor[rgb]{0.13,0.29,0.53}{#1}}
\newcommand{\BaseNTok}[1]{\textcolor[rgb]{0.00,0.00,0.81}{#1}}
\newcommand{\BuiltInTok}[1]{#1}
\newcommand{\CharTok}[1]{\textcolor[rgb]{0.31,0.60,0.02}{#1}}
\newcommand{\CommentTok}[1]{\textcolor[rgb]{0.56,0.35,0.01}{\textit{#1}}}
\newcommand{\CommentVarTok}[1]{\textcolor[rgb]{0.56,0.35,0.01}{\textbf{\textit{#1}}}}
\newcommand{\ConstantTok}[1]{\textcolor[rgb]{0.56,0.35,0.01}{#1}}
\newcommand{\ControlFlowTok}[1]{\textcolor[rgb]{0.13,0.29,0.53}{\textbf{#1}}}
\newcommand{\DataTypeTok}[1]{\textcolor[rgb]{0.13,0.29,0.53}{#1}}
\newcommand{\DecValTok}[1]{\textcolor[rgb]{0.00,0.00,0.81}{#1}}
\newcommand{\DocumentationTok}[1]{\textcolor[rgb]{0.56,0.35,0.01}{\textbf{\textit{#1}}}}
\newcommand{\ErrorTok}[1]{\textcolor[rgb]{0.64,0.00,0.00}{\textbf{#1}}}
\newcommand{\ExtensionTok}[1]{#1}
\newcommand{\FloatTok}[1]{\textcolor[rgb]{0.00,0.00,0.81}{#1}}
\newcommand{\FunctionTok}[1]{\textcolor[rgb]{0.13,0.29,0.53}{\textbf{#1}}}
\newcommand{\ImportTok}[1]{#1}
\newcommand{\InformationTok}[1]{\textcolor[rgb]{0.56,0.35,0.01}{\textbf{\textit{#1}}}}
\newcommand{\KeywordTok}[1]{\textcolor[rgb]{0.13,0.29,0.53}{\textbf{#1}}}
\newcommand{\NormalTok}[1]{#1}
\newcommand{\OperatorTok}[1]{\textcolor[rgb]{0.81,0.36,0.00}{\textbf{#1}}}
\newcommand{\OtherTok}[1]{\textcolor[rgb]{0.56,0.35,0.01}{#1}}
\newcommand{\PreprocessorTok}[1]{\textcolor[rgb]{0.56,0.35,0.01}{\textit{#1}}}
\newcommand{\RegionMarkerTok}[1]{#1}
\newcommand{\SpecialCharTok}[1]{\textcolor[rgb]{0.81,0.36,0.00}{\textbf{#1}}}
\newcommand{\SpecialStringTok}[1]{\textcolor[rgb]{0.31,0.60,0.02}{#1}}
\newcommand{\StringTok}[1]{\textcolor[rgb]{0.31,0.60,0.02}{#1}}
\newcommand{\VariableTok}[1]{\textcolor[rgb]{0.00,0.00,0.00}{#1}}
\newcommand{\VerbatimStringTok}[1]{\textcolor[rgb]{0.31,0.60,0.02}{#1}}
\newcommand{\WarningTok}[1]{\textcolor[rgb]{0.56,0.35,0.01}{\textbf{\textit{#1}}}}
\usepackage{graphicx}
\makeatletter
\newsavebox\pandoc@box
\newcommand*\pandocbounded[1]{% scales image to fit in text height/width
  \sbox\pandoc@box{#1}%
  \Gscale@div\@tempa{\textheight}{\dimexpr\ht\pandoc@box+\dp\pandoc@box\relax}%
  \Gscale@div\@tempb{\linewidth}{\wd\pandoc@box}%
  \ifdim\@tempb\p@<\@tempa\p@\let\@tempa\@tempb\fi% select the smaller of both
  \ifdim\@tempa\p@<\p@\scalebox{\@tempa}{\usebox\pandoc@box}%
  \else\usebox{\pandoc@box}%
  \fi%
}
% Set default figure placement to htbp
\def\fps@figure{htbp}
\makeatother
\setlength{\emergencystretch}{3em} % prevent overfull lines
\providecommand{\tightlist}{%
  \setlength{\itemsep}{0pt}\setlength{\parskip}{0pt}}
\usepackage{bookmark}
\IfFileExists{xurl.sty}{\usepackage{xurl}}{} % add URL line breaks if available
\urlstyle{same}
\hypersetup{
  pdftitle={STAT 345 Midterm Project},
  hidelinks,
  pdfcreator={LaTeX via pandoc}}

\title{STAT 345 Midterm Project}
\author{}
\date{\vspace{-2.5em}April 3rd, 2026}

\begin{document}
\maketitle

\begin{quote}
``Our offense is like the Pythagorean Theorem. There is no answer!'',
\emph{Shaquille O'Neal}
\end{quote}

\subsection{The Background}\label{the-background}

Your role for the midterm project is that of data analyst intern at an
NBA (professional basketball) team. Your direct supervisor (also part of
the analytics team) has asked you to create a data visualization to
illustrate how (or if) the team's shots have changed over time. After
some initial clarifying questions, your supervisor confessed that they
had seen some pretty cool shot charts at
\url{http://savvastjortjoglou.com/nba-shot-sharts.html} and would like
to extend the ideas a bit.

Your data for the midterm project may come from a variety of sources,
including the NBA directly, as well as Basketball-Reference, HoopsHype,
and others. There are several ways to access the data, but perhaps the
simplest is through a package. There are a few options (each with some
benefits, drawbacks, and potential headaches):

\begin{itemize}
\tightlist
\item
  \texttt{nbastatR}
\item
  \texttt{hoopR}
\item
  \texttt{AdvancedBasketballStats}
\end{itemize}

Note: if you would prefer to analyze college basketball, packages exist
for that too (such as \texttt{cbbdata}). If you would prefer to look
into a different sport, we can investigate that as well. Please let me
know in advance if you intend to go down one of these alternate routes.

\subsection{The Tasks}\label{the-tasks}

\begin{enumerate}
\def\labelenumi{\arabic{enumi}.}
\item
  (30 points) Produce a graphic displaying the shot locations for a
  particular team over several years. Some notes:

  \begin{itemize}
  \tightlist
  \item
    Colors should be chosen to reflect the team, if possible.
  \item
    There are likely many overlaid points -- handle this by either
    binning these by location, or use opacity and size.
  \item
    Incorporate information about whether the shot was made or not
    (shape, color, etc.).
  \item
    The graphic should be well-labeled, titled, etc.
  \item
    Start with a graph for a single year, then extend to several years.
    Up to 20 years of shot data is available. Either facet these by year
    or animate using the years.
  \item
    You'll want to figure out what the coordinates mean somehow. This
    might be through the documentation, but could also be determined
    using aspects of the data itself and the dimensions of an NBA court.
  \item
    Put a basketball court on the background of the image (you'll need
    to scale it appropriately).
  \end{itemize}
\end{enumerate}

\begin{verbatim}
## -- Attaching core tidyverse packages ------------------------ tidyverse 2.0.0 --
## v dplyr     1.1.4     v readr     2.1.6
## v forcats   1.0.1     v stringr   1.6.0
## v ggplot2   4.0.1     v tibble    3.3.1
## v lubridate 1.9.4     v tidyr     1.3.2
## v purrr     1.2.1     
## -- Conflicts ------------------------------------------ tidyverse_conflicts() --
## x dplyr::filter() masks stats::filter()
## x dplyr::lag()    masks stats::lag()
## i Use the conflicted package (<http://conflicted.r-lib.org/>) to force all conflicts to become errors
\end{verbatim}

\begin{verbatim}
## Warning: package 'ggpubr' was built under R version 4.5.3
\end{verbatim}

\begin{verbatim}
## Warning: package 'png' was built under R version 4.5.3
\end{verbatim}

\begin{verbatim}
## Warning: package 'gganimate' was built under R version 4.5.3
\end{verbatim}

\begin{verbatim}
## Warning: package 'gifski' was built under R version 4.5.3
\end{verbatim}

\begin{verbatim}
## -- ESPN NBA Play-by-Play from hoopR data repository ------------- hoopR 2.1.0 --
\end{verbatim}

\begin{verbatim}
## i Data updated: 2026-04-02 06:44:12 CDT
\end{verbatim}

\begin{verbatim}
## # A tibble: 108,641 x 7
##    score_value team_id shooting_play coordinate_x_raw coordinate_y_raw season
##          <int>   <int> <lgl>                    <dbl>            <dbl>  <int>
##  1           0       3 TRUE                        23              1     2015
##  2           2       3 TRUE                        20             19     2015
##  3           2       3 TRUE                        24             15     2015
##  4           2       3 TRUE                        24             17     2015
##  5           2       3 TRUE                        28              2     2015
##  6           0       3 TRUE                        25             13.8   2015
##  7           0       3 TRUE                        26              1     2015
##  8           0       3 TRUE                        35             23     2015
##  9           2       3 TRUE                        12              2     2015
## 10           0       3 TRUE                        22              5     2015
## # i 108,631 more rows
## # i 1 more variable: make_miss <chr>
\end{verbatim}

\begin{verbatim}
## Warning in formals(fun): argument is not a function
\end{verbatim}

\pandocbounded{\includegraphics[keepaspectratio]{midterm-1_files/figure-latex/unnamed-chunk-2-1.pdf}}
\pandocbounded{\includegraphics[keepaspectratio]{midterm-1_files/figure-latex/unnamed-chunk-2-2.pdf}}
\pandocbounded{\includegraphics[keepaspectratio]{midterm-1_files/figure-latex/unnamed-chunk-2-3.pdf}}
\pandocbounded{\includegraphics[keepaspectratio]{midterm-1_files/figure-latex/unnamed-chunk-2-4.pdf}}
\pandocbounded{\includegraphics[keepaspectratio]{midterm-1_files/figure-latex/unnamed-chunk-2-5.pdf}}
\pandocbounded{\includegraphics[keepaspectratio]{midterm-1_files/figure-latex/unnamed-chunk-2-6.pdf}}
\pandocbounded{\includegraphics[keepaspectratio]{midterm-1_files/figure-latex/unnamed-chunk-2-7.pdf}}
\pandocbounded{\includegraphics[keepaspectratio]{midterm-1_files/figure-latex/unnamed-chunk-2-8.pdf}}
\pandocbounded{\includegraphics[keepaspectratio]{midterm-1_files/figure-latex/unnamed-chunk-2-9.pdf}}
\pandocbounded{\includegraphics[keepaspectratio]{midterm-1_files/figure-latex/unnamed-chunk-2-10.pdf}}
\pandocbounded{\includegraphics[keepaspectratio]{midterm-1_files/figure-latex/unnamed-chunk-2-11.pdf}}
\pandocbounded{\includegraphics[keepaspectratio]{midterm-1_files/figure-latex/unnamed-chunk-2-12.pdf}}
\pandocbounded{\includegraphics[keepaspectratio]{midterm-1_files/figure-latex/unnamed-chunk-2-13.pdf}}
\pandocbounded{\includegraphics[keepaspectratio]{midterm-1_files/figure-latex/unnamed-chunk-2-14.pdf}}
\pandocbounded{\includegraphics[keepaspectratio]{midterm-1_files/figure-latex/unnamed-chunk-2-15.pdf}}
\pandocbounded{\includegraphics[keepaspectratio]{midterm-1_files/figure-latex/unnamed-chunk-2-16.pdf}}
\pandocbounded{\includegraphics[keepaspectratio]{midterm-1_files/figure-latex/unnamed-chunk-2-17.pdf}}
\pandocbounded{\includegraphics[keepaspectratio]{midterm-1_files/figure-latex/unnamed-chunk-2-18.pdf}}
\pandocbounded{\includegraphics[keepaspectratio]{midterm-1_files/figure-latex/unnamed-chunk-2-19.pdf}}
\pandocbounded{\includegraphics[keepaspectratio]{midterm-1_files/figure-latex/unnamed-chunk-2-20.pdf}}
\pandocbounded{\includegraphics[keepaspectratio]{midterm-1_files/figure-latex/unnamed-chunk-2-21.pdf}}
\pandocbounded{\includegraphics[keepaspectratio]{midterm-1_files/figure-latex/unnamed-chunk-2-22.pdf}}
\pandocbounded{\includegraphics[keepaspectratio]{midterm-1_files/figure-latex/unnamed-chunk-2-23.pdf}}
\pandocbounded{\includegraphics[keepaspectratio]{midterm-1_files/figure-latex/unnamed-chunk-2-24.pdf}}
\pandocbounded{\includegraphics[keepaspectratio]{midterm-1_files/figure-latex/unnamed-chunk-2-25.pdf}}
\pandocbounded{\includegraphics[keepaspectratio]{midterm-1_files/figure-latex/unnamed-chunk-2-26.pdf}}
\pandocbounded{\includegraphics[keepaspectratio]{midterm-1_files/figure-latex/unnamed-chunk-2-27.pdf}}
\pandocbounded{\includegraphics[keepaspectratio]{midterm-1_files/figure-latex/unnamed-chunk-2-28.pdf}}
\pandocbounded{\includegraphics[keepaspectratio]{midterm-1_files/figure-latex/unnamed-chunk-2-29.pdf}}
\pandocbounded{\includegraphics[keepaspectratio]{midterm-1_files/figure-latex/unnamed-chunk-2-30.pdf}}
\pandocbounded{\includegraphics[keepaspectratio]{midterm-1_files/figure-latex/unnamed-chunk-2-31.pdf}}
\pandocbounded{\includegraphics[keepaspectratio]{midterm-1_files/figure-latex/unnamed-chunk-2-32.pdf}}
\pandocbounded{\includegraphics[keepaspectratio]{midterm-1_files/figure-latex/unnamed-chunk-2-33.pdf}}
\pandocbounded{\includegraphics[keepaspectratio]{midterm-1_files/figure-latex/unnamed-chunk-2-34.pdf}}
\pandocbounded{\includegraphics[keepaspectratio]{midterm-1_files/figure-latex/unnamed-chunk-2-35.pdf}}
\pandocbounded{\includegraphics[keepaspectratio]{midterm-1_files/figure-latex/unnamed-chunk-2-36.pdf}}
\pandocbounded{\includegraphics[keepaspectratio]{midterm-1_files/figure-latex/unnamed-chunk-2-37.pdf}}
\pandocbounded{\includegraphics[keepaspectratio]{midterm-1_files/figure-latex/unnamed-chunk-2-38.pdf}}
\pandocbounded{\includegraphics[keepaspectratio]{midterm-1_files/figure-latex/unnamed-chunk-2-39.pdf}}
\pandocbounded{\includegraphics[keepaspectratio]{midterm-1_files/figure-latex/unnamed-chunk-2-40.pdf}}
\pandocbounded{\includegraphics[keepaspectratio]{midterm-1_files/figure-latex/unnamed-chunk-2-41.pdf}}
\pandocbounded{\includegraphics[keepaspectratio]{midterm-1_files/figure-latex/unnamed-chunk-2-42.pdf}}
\pandocbounded{\includegraphics[keepaspectratio]{midterm-1_files/figure-latex/unnamed-chunk-2-43.pdf}}
\pandocbounded{\includegraphics[keepaspectratio]{midterm-1_files/figure-latex/unnamed-chunk-2-44.pdf}}
\pandocbounded{\includegraphics[keepaspectratio]{midterm-1_files/figure-latex/unnamed-chunk-2-45.pdf}}
\pandocbounded{\includegraphics[keepaspectratio]{midterm-1_files/figure-latex/unnamed-chunk-2-46.pdf}}
\pandocbounded{\includegraphics[keepaspectratio]{midterm-1_files/figure-latex/unnamed-chunk-2-47.pdf}}
\pandocbounded{\includegraphics[keepaspectratio]{midterm-1_files/figure-latex/unnamed-chunk-2-48.pdf}}
\pandocbounded{\includegraphics[keepaspectratio]{midterm-1_files/figure-latex/unnamed-chunk-2-49.pdf}}
\pandocbounded{\includegraphics[keepaspectratio]{midterm-1_files/figure-latex/unnamed-chunk-2-50.pdf}}
\pandocbounded{\includegraphics[keepaspectratio]{midterm-1_files/figure-latex/unnamed-chunk-2-51.pdf}}
\pandocbounded{\includegraphics[keepaspectratio]{midterm-1_files/figure-latex/unnamed-chunk-2-52.pdf}}
\pandocbounded{\includegraphics[keepaspectratio]{midterm-1_files/figure-latex/unnamed-chunk-2-53.pdf}}
\pandocbounded{\includegraphics[keepaspectratio]{midterm-1_files/figure-latex/unnamed-chunk-2-54.pdf}}
\pandocbounded{\includegraphics[keepaspectratio]{midterm-1_files/figure-latex/unnamed-chunk-2-55.pdf}}
\pandocbounded{\includegraphics[keepaspectratio]{midterm-1_files/figure-latex/unnamed-chunk-2-56.pdf}}
\pandocbounded{\includegraphics[keepaspectratio]{midterm-1_files/figure-latex/unnamed-chunk-2-57.pdf}}
\pandocbounded{\includegraphics[keepaspectratio]{midterm-1_files/figure-latex/unnamed-chunk-2-58.pdf}}
\pandocbounded{\includegraphics[keepaspectratio]{midterm-1_files/figure-latex/unnamed-chunk-2-59.pdf}}
\pandocbounded{\includegraphics[keepaspectratio]{midterm-1_files/figure-latex/unnamed-chunk-2-60.pdf}}
\pandocbounded{\includegraphics[keepaspectratio]{midterm-1_files/figure-latex/unnamed-chunk-2-61.pdf}}
\pandocbounded{\includegraphics[keepaspectratio]{midterm-1_files/figure-latex/unnamed-chunk-2-62.pdf}}
\pandocbounded{\includegraphics[keepaspectratio]{midterm-1_files/figure-latex/unnamed-chunk-2-63.pdf}}
\pandocbounded{\includegraphics[keepaspectratio]{midterm-1_files/figure-latex/unnamed-chunk-2-64.pdf}}
\pandocbounded{\includegraphics[keepaspectratio]{midterm-1_files/figure-latex/unnamed-chunk-2-65.pdf}}
\pandocbounded{\includegraphics[keepaspectratio]{midterm-1_files/figure-latex/unnamed-chunk-2-66.pdf}}
\pandocbounded{\includegraphics[keepaspectratio]{midterm-1_files/figure-latex/unnamed-chunk-2-67.pdf}}
\pandocbounded{\includegraphics[keepaspectratio]{midterm-1_files/figure-latex/unnamed-chunk-2-68.pdf}}
\pandocbounded{\includegraphics[keepaspectratio]{midterm-1_files/figure-latex/unnamed-chunk-2-69.pdf}}
\pandocbounded{\includegraphics[keepaspectratio]{midterm-1_files/figure-latex/unnamed-chunk-2-70.pdf}}
\pandocbounded{\includegraphics[keepaspectratio]{midterm-1_files/figure-latex/unnamed-chunk-2-71.pdf}}
\pandocbounded{\includegraphics[keepaspectratio]{midterm-1_files/figure-latex/unnamed-chunk-2-72.pdf}}
\pandocbounded{\includegraphics[keepaspectratio]{midterm-1_files/figure-latex/unnamed-chunk-2-73.pdf}}
\pandocbounded{\includegraphics[keepaspectratio]{midterm-1_files/figure-latex/unnamed-chunk-2-74.pdf}}
\pandocbounded{\includegraphics[keepaspectratio]{midterm-1_files/figure-latex/unnamed-chunk-2-75.pdf}}
\pandocbounded{\includegraphics[keepaspectratio]{midterm-1_files/figure-latex/unnamed-chunk-2-76.pdf}}
\pandocbounded{\includegraphics[keepaspectratio]{midterm-1_files/figure-latex/unnamed-chunk-2-77.pdf}}
\pandocbounded{\includegraphics[keepaspectratio]{midterm-1_files/figure-latex/unnamed-chunk-2-78.pdf}}
\pandocbounded{\includegraphics[keepaspectratio]{midterm-1_files/figure-latex/unnamed-chunk-2-79.pdf}}
\pandocbounded{\includegraphics[keepaspectratio]{midterm-1_files/figure-latex/unnamed-chunk-2-80.pdf}}
\pandocbounded{\includegraphics[keepaspectratio]{midterm-1_files/figure-latex/unnamed-chunk-2-81.pdf}}
\pandocbounded{\includegraphics[keepaspectratio]{midterm-1_files/figure-latex/unnamed-chunk-2-82.pdf}}
\pandocbounded{\includegraphics[keepaspectratio]{midterm-1_files/figure-latex/unnamed-chunk-2-83.pdf}}
\pandocbounded{\includegraphics[keepaspectratio]{midterm-1_files/figure-latex/unnamed-chunk-2-84.pdf}}
\pandocbounded{\includegraphics[keepaspectratio]{midterm-1_files/figure-latex/unnamed-chunk-2-85.pdf}}
\pandocbounded{\includegraphics[keepaspectratio]{midterm-1_files/figure-latex/unnamed-chunk-2-86.pdf}}
\pandocbounded{\includegraphics[keepaspectratio]{midterm-1_files/figure-latex/unnamed-chunk-2-87.pdf}}
\pandocbounded{\includegraphics[keepaspectratio]{midterm-1_files/figure-latex/unnamed-chunk-2-88.pdf}}
\pandocbounded{\includegraphics[keepaspectratio]{midterm-1_files/figure-latex/unnamed-chunk-2-89.pdf}}
\pandocbounded{\includegraphics[keepaspectratio]{midterm-1_files/figure-latex/unnamed-chunk-2-90.pdf}}
\pandocbounded{\includegraphics[keepaspectratio]{midterm-1_files/figure-latex/unnamed-chunk-2-91.pdf}}
\pandocbounded{\includegraphics[keepaspectratio]{midterm-1_files/figure-latex/unnamed-chunk-2-92.pdf}}
\pandocbounded{\includegraphics[keepaspectratio]{midterm-1_files/figure-latex/unnamed-chunk-2-93.pdf}}
\pandocbounded{\includegraphics[keepaspectratio]{midterm-1_files/figure-latex/unnamed-chunk-2-94.pdf}}
\pandocbounded{\includegraphics[keepaspectratio]{midterm-1_files/figure-latex/unnamed-chunk-2-95.pdf}}
\pandocbounded{\includegraphics[keepaspectratio]{midterm-1_files/figure-latex/unnamed-chunk-2-96.pdf}}
\pandocbounded{\includegraphics[keepaspectratio]{midterm-1_files/figure-latex/unnamed-chunk-2-97.pdf}}
\pandocbounded{\includegraphics[keepaspectratio]{midterm-1_files/figure-latex/unnamed-chunk-2-98.pdf}}
\pandocbounded{\includegraphics[keepaspectratio]{midterm-1_files/figure-latex/unnamed-chunk-2-99.pdf}}
\pandocbounded{\includegraphics[keepaspectratio]{midterm-1_files/figure-latex/unnamed-chunk-2-100.pdf}}

\begin{enumerate}
\def\labelenumi{\arabic{enumi}.}
\setcounter{enumi}{1}
\item
  (30 points) Summarize the graphic/series of graphics into a
  digestible, bullet-point brief report for front-office staff. Some
  notes:

  \begin{itemize}
  \tightlist
  \item
    The main body of the report should be very brief -- just the
    graphic(s) and the bullet-pointed list of findings, which should be
    short and clear.
  \item
    Include a more detailed explanation of these bullet points, for
    further reading by those interested. This section should follow the
    bullet-point section, but should be organized similarly for
    reference.
  \item
    Your report to the front-office shouldn't include any code.
  \item
    This report should be generated using RMarkdown. However, the choice
    of output type (Word, PDF, or HTML) is up to you (you could even
    make slides if you want to).
  \end{itemize}
\end{enumerate}

\subsubsection{\texorpdfstring{\textbf{Summary:}}{Summary:}}\label{summary}

\begin{itemize}
\item
  \textbf{Most Common Shot Locations:} under the basket, free throw
  line, corners, and top of three point line
\item
  \textbf{Shift Across Seasons:} shot distribution slowly shifts from a
  wide spread to a condensed shot selection of just a few main spots
\item
  \textbf{Make/Miss Ratio:} made shots appear more in the common shot
  locations and increase over time
\item
  \textbf{Suggestions for Improvement:} focus on improving shot
  performance or shot location to increase overall scoring success
\end{itemize}

\subsubsection{\texorpdfstring{\textbf{Explanation:}}{Explanation:}}\label{explanation}

\begin{itemize}
\item
  \textbf{Most Common Shot Locations:}

  \begin{itemize}
  \tightlist
  \item
    The majority of shot attempts are made directly under the basket, on
    the free throw line, or along the three point line. On the three
    point line, there is a preference for the corners and wings, with
    fewer in between. The highest concentration of shot attempts are
    under the basket. This aligns with the common strategy for most
    basketball teams, which is to focus on point efficiency by
    prioritizing either close range shots and layups for easy points or
    three point shots for more points.
  \end{itemize}
\item
  \textbf{Shift Across Seasons:}

  \begin{itemize}
  \tightlist
  \item
    Looking at all the seasons, from 2015 to 2026, shot distribution
    gradually shifts. In the early seasons, there is a large spread of
    shots all over the court. Then, over time, the shot distribution
    narrows down to most being taken under the basket, on the free throw
    line, in the corners, or on the wings and less being taken in the
    mid range area. This indicates a clear improvement in playing style
    to prioritize scoring efficiency.
  \end{itemize}
\item
  \textbf{Make/Miss Ratio:}

  \begin{itemize}
  \tightlist
  \item
    The majority of made shots tend to be concentrated under the basket,
    where scoring is easier and more consistent, as well as on the free
    throw line in later seasons, where shots are uncontested. Over time,
    the data shows improvement in shot successes. In the early seasons
    you can clearly see there are more red dots, especially for free
    throws. It slowly shifts from red to blue in the later seasons ,
    showing an increase in made free throws. There also appears to be a
    greater percentage of made shots beneath the basket in the later
    seasons compared to the early seasons.
  \end{itemize}
\item
  \textbf{Suggestions for Improvement:}

  \begin{itemize}
  \tightlist
  \item
    To improve overall scoring success, I suggest focusing on shot
    performance and better ball movement. Quality shots and opening up
    players would help increase the percentage of made shots. Another
    suggestion is to focus on taking more close range shots rather than
    three-pointers, since there is already higher success under the
    basket, which could lead to higher scoring output.
  \end{itemize}
\end{itemize}

\begin{enumerate}
\def\labelenumi{\arabic{enumi}.}
\setcounter{enumi}{2}
\tightlist
\item
  (30 points) Write and document clean, efficient, reproducible code.
  Some notes:
\end{enumerate}

\begin{itemize}
\item
  This code will be viewed by your direct supervisor.
\item
  The code file should include your code to gather, join, and clean the
  data; the code to generate the graphic(s) presented; and your
  commentary on the results (so, a single .rmd file, or an .rmd file
  that sources an .r file).
\item
  Your code should be clean, organized, and reproducible. Remove
  unnecessary/scratch/exploratory code.
\item
  Your code should be well commented. In particular, any decisions or
  judgement calls made in the analysis process should be
  explained/justified. Sections of code should be identified even if not
  functionalized (including purpose, data/argument inputs, analysis
  outputs).
\end{itemize}

\begin{Shaded}
\begin{Highlighting}[]
\CommentTok{\# add libraries}
\FunctionTok{library}\NormalTok{(tidyverse)}
\FunctionTok{library}\NormalTok{(hoopR)}
\FunctionTok{library}\NormalTok{(dplyr)}
\FunctionTok{library}\NormalTok{(ggplot2)}
\FunctionTok{library}\NormalTok{(knitr)}
\FunctionTok{library}\NormalTok{(ggpubr)}
\FunctionTok{library}\NormalTok{(png)}
\FunctionTok{library}\NormalTok{(gganimate)}
\FunctionTok{library}\NormalTok{(gifski)}

\CommentTok{\# load in data from hoopR for seasons 2015 to 2026}
\NormalTok{data }\OtherTok{\textless{}{-}}\NormalTok{ hoopR}\SpecialCharTok{::}\FunctionTok{load\_nba\_pbp}\NormalTok{(}\AttributeTok{seasons =} \DecValTok{2015}\SpecialCharTok{:}\DecValTok{2026}\NormalTok{)}

\CommentTok{\# filters data to the New Orleans Pelicans (team 3) and only shooting plays}
\CommentTok{\# adds new column to classify makes/misses}
\CommentTok{\# gets rid of unnecessary columns so data is easier to manage and loads faster}
\NormalTok{baby\_data }\OtherTok{\textless{}{-}}\NormalTok{ data }\SpecialCharTok{|\textgreater{}} 
  \FunctionTok{filter}\NormalTok{(team\_id }\SpecialCharTok{==} \DecValTok{3}\NormalTok{, shooting\_play }\SpecialCharTok{==} \ConstantTok{TRUE}\NormalTok{) }\SpecialCharTok{|\textgreater{}} 
  \FunctionTok{mutate}\NormalTok{(}\AttributeTok{make\_miss =} \FunctionTok{ifelse}\NormalTok{(score\_value }\SpecialCharTok{\textgreater{}} \DecValTok{0}\NormalTok{, }\StringTok{"Make"}\NormalTok{, }\StringTok{"Miss"}\NormalTok{)) }\SpecialCharTok{|\textgreater{}} 
  \FunctionTok{select}\NormalTok{(}\StringTok{"score\_value"}\NormalTok{, }\StringTok{"team\_id"}\NormalTok{, }\StringTok{"shooting\_play"}\NormalTok{, }\StringTok{"coordinate\_x\_raw"}\NormalTok{, }\StringTok{"coordinate\_y\_raw"}\NormalTok{, }\StringTok{"season"}\NormalTok{, }\StringTok{"make\_miss"}\NormalTok{)}

\CommentTok{\# gets background image of court for animation}
\NormalTok{img }\OtherTok{\textless{}{-}} \FunctionTok{readPNG}\NormalTok{(}\StringTok{"midterm\_bbcourt.png"}\NormalTok{)}

\CommentTok{\# creates a plot based on raw x and y coords}
\CommentTok{\# adds background image of NBA court}
\CommentTok{\# adjust opacity (alpha) so darker AND added weight/height so dots are larger for multiple shots from same coords}
\CommentTok{\# assign colors for made/missed shots...made shots = Pelicans Navy, missed = Pelicans Red}
\CommentTok{\# added labels for title, axises, and legend}
\CommentTok{\# fix coords for more accurate alignment with background image}
\CommentTok{\# add animation between seasons, adjusting transition/state length for better viewing}
\FunctionTok{ggplot}\NormalTok{(baby\_data, }\FunctionTok{aes}\NormalTok{(}\AttributeTok{x =}\NormalTok{ coordinate\_x\_raw, }\AttributeTok{y =}\NormalTok{ coordinate\_y\_raw)) }\SpecialCharTok{+} 
  \FunctionTok{background\_image}\NormalTok{(img) }\SpecialCharTok{+} 
  \FunctionTok{geom\_jitter}\NormalTok{(}\FunctionTok{aes}\NormalTok{(}\AttributeTok{color =}\NormalTok{ make\_miss), }\AttributeTok{width =} \FloatTok{0.5}\NormalTok{, }\AttributeTok{height =} \FloatTok{0.5}\NormalTok{, }\AttributeTok{alpha =} \FloatTok{0.1}\NormalTok{, }\AttributeTok{size =} \DecValTok{1}\NormalTok{) }\SpecialCharTok{+} 
  \FunctionTok{scale\_color\_manual}\NormalTok{(}\AttributeTok{values =} \FunctionTok{c}\NormalTok{(}\StringTok{"Make"} \OtherTok{=} \StringTok{"\#0C2340"}\NormalTok{, }\StringTok{"Miss"} \OtherTok{=} \StringTok{"\#C8102E"}\NormalTok{)) }\SpecialCharTok{+} 
  \FunctionTok{labs}\NormalTok{(}\AttributeTok{title =} \StringTok{"New Orleans Pelicans Shots \{closest\_state\}"}\NormalTok{, }\AttributeTok{x =} \StringTok{"Raw X Coord"}\NormalTok{, }\AttributeTok{y =} \StringTok{"Raw Y Coord"}\NormalTok{, }\AttributeTok{color =} \StringTok{"Shot Result"}\NormalTok{) }\SpecialCharTok{+} 
  \FunctionTok{coord\_fixed}\NormalTok{(}\AttributeTok{xlim =} \FunctionTok{c}\NormalTok{(}\DecValTok{0}\NormalTok{, }\DecValTok{50}\NormalTok{), }\AttributeTok{ylim =} \FunctionTok{c}\NormalTok{(}\DecValTok{0}\NormalTok{, }\FloatTok{36.5}\NormalTok{)) }\SpecialCharTok{+} 
  \FunctionTok{transition\_states}\NormalTok{(season, }\AttributeTok{transition\_length =} \DecValTok{3}\NormalTok{, }\AttributeTok{state\_length =} \DecValTok{100}\NormalTok{)}
\end{Highlighting}
\end{Shaded}

\begin{verbatim}
## Warning in formals(fun): argument is not a function
\end{verbatim}

\pandocbounded{\includegraphics[keepaspectratio]{midterm-1_files/figure-latex/unnamed-chunk-3-1.pdf}}
\pandocbounded{\includegraphics[keepaspectratio]{midterm-1_files/figure-latex/unnamed-chunk-3-2.pdf}}
\pandocbounded{\includegraphics[keepaspectratio]{midterm-1_files/figure-latex/unnamed-chunk-3-3.pdf}}
\pandocbounded{\includegraphics[keepaspectratio]{midterm-1_files/figure-latex/unnamed-chunk-3-4.pdf}}
\pandocbounded{\includegraphics[keepaspectratio]{midterm-1_files/figure-latex/unnamed-chunk-3-5.pdf}}
\pandocbounded{\includegraphics[keepaspectratio]{midterm-1_files/figure-latex/unnamed-chunk-3-6.pdf}}
\pandocbounded{\includegraphics[keepaspectratio]{midterm-1_files/figure-latex/unnamed-chunk-3-7.pdf}}
\pandocbounded{\includegraphics[keepaspectratio]{midterm-1_files/figure-latex/unnamed-chunk-3-8.pdf}}
\pandocbounded{\includegraphics[keepaspectratio]{midterm-1_files/figure-latex/unnamed-chunk-3-9.pdf}}
\pandocbounded{\includegraphics[keepaspectratio]{midterm-1_files/figure-latex/unnamed-chunk-3-10.pdf}}
\pandocbounded{\includegraphics[keepaspectratio]{midterm-1_files/figure-latex/unnamed-chunk-3-11.pdf}}
\pandocbounded{\includegraphics[keepaspectratio]{midterm-1_files/figure-latex/unnamed-chunk-3-12.pdf}}
\pandocbounded{\includegraphics[keepaspectratio]{midterm-1_files/figure-latex/unnamed-chunk-3-13.pdf}}
\pandocbounded{\includegraphics[keepaspectratio]{midterm-1_files/figure-latex/unnamed-chunk-3-14.pdf}}
\pandocbounded{\includegraphics[keepaspectratio]{midterm-1_files/figure-latex/unnamed-chunk-3-15.pdf}}
\pandocbounded{\includegraphics[keepaspectratio]{midterm-1_files/figure-latex/unnamed-chunk-3-16.pdf}}
\pandocbounded{\includegraphics[keepaspectratio]{midterm-1_files/figure-latex/unnamed-chunk-3-17.pdf}}
\pandocbounded{\includegraphics[keepaspectratio]{midterm-1_files/figure-latex/unnamed-chunk-3-18.pdf}}
\pandocbounded{\includegraphics[keepaspectratio]{midterm-1_files/figure-latex/unnamed-chunk-3-19.pdf}}
\pandocbounded{\includegraphics[keepaspectratio]{midterm-1_files/figure-latex/unnamed-chunk-3-20.pdf}}
\pandocbounded{\includegraphics[keepaspectratio]{midterm-1_files/figure-latex/unnamed-chunk-3-21.pdf}}
\pandocbounded{\includegraphics[keepaspectratio]{midterm-1_files/figure-latex/unnamed-chunk-3-22.pdf}}
\pandocbounded{\includegraphics[keepaspectratio]{midterm-1_files/figure-latex/unnamed-chunk-3-23.pdf}}
\pandocbounded{\includegraphics[keepaspectratio]{midterm-1_files/figure-latex/unnamed-chunk-3-24.pdf}}
\pandocbounded{\includegraphics[keepaspectratio]{midterm-1_files/figure-latex/unnamed-chunk-3-25.pdf}}
\pandocbounded{\includegraphics[keepaspectratio]{midterm-1_files/figure-latex/unnamed-chunk-3-26.pdf}}
\pandocbounded{\includegraphics[keepaspectratio]{midterm-1_files/figure-latex/unnamed-chunk-3-27.pdf}}
\pandocbounded{\includegraphics[keepaspectratio]{midterm-1_files/figure-latex/unnamed-chunk-3-28.pdf}}
\pandocbounded{\includegraphics[keepaspectratio]{midterm-1_files/figure-latex/unnamed-chunk-3-29.pdf}}
\pandocbounded{\includegraphics[keepaspectratio]{midterm-1_files/figure-latex/unnamed-chunk-3-30.pdf}}
\pandocbounded{\includegraphics[keepaspectratio]{midterm-1_files/figure-latex/unnamed-chunk-3-31.pdf}}
\pandocbounded{\includegraphics[keepaspectratio]{midterm-1_files/figure-latex/unnamed-chunk-3-32.pdf}}
\pandocbounded{\includegraphics[keepaspectratio]{midterm-1_files/figure-latex/unnamed-chunk-3-33.pdf}}
\pandocbounded{\includegraphics[keepaspectratio]{midterm-1_files/figure-latex/unnamed-chunk-3-34.pdf}}
\pandocbounded{\includegraphics[keepaspectratio]{midterm-1_files/figure-latex/unnamed-chunk-3-35.pdf}}
\pandocbounded{\includegraphics[keepaspectratio]{midterm-1_files/figure-latex/unnamed-chunk-3-36.pdf}}
\pandocbounded{\includegraphics[keepaspectratio]{midterm-1_files/figure-latex/unnamed-chunk-3-37.pdf}}
\pandocbounded{\includegraphics[keepaspectratio]{midterm-1_files/figure-latex/unnamed-chunk-3-38.pdf}}
\pandocbounded{\includegraphics[keepaspectratio]{midterm-1_files/figure-latex/unnamed-chunk-3-39.pdf}}
\pandocbounded{\includegraphics[keepaspectratio]{midterm-1_files/figure-latex/unnamed-chunk-3-40.pdf}}
\pandocbounded{\includegraphics[keepaspectratio]{midterm-1_files/figure-latex/unnamed-chunk-3-41.pdf}}
\pandocbounded{\includegraphics[keepaspectratio]{midterm-1_files/figure-latex/unnamed-chunk-3-42.pdf}}
\pandocbounded{\includegraphics[keepaspectratio]{midterm-1_files/figure-latex/unnamed-chunk-3-43.pdf}}
\pandocbounded{\includegraphics[keepaspectratio]{midterm-1_files/figure-latex/unnamed-chunk-3-44.pdf}}
\pandocbounded{\includegraphics[keepaspectratio]{midterm-1_files/figure-latex/unnamed-chunk-3-45.pdf}}
\pandocbounded{\includegraphics[keepaspectratio]{midterm-1_files/figure-latex/unnamed-chunk-3-46.pdf}}
\pandocbounded{\includegraphics[keepaspectratio]{midterm-1_files/figure-latex/unnamed-chunk-3-47.pdf}}
\pandocbounded{\includegraphics[keepaspectratio]{midterm-1_files/figure-latex/unnamed-chunk-3-48.pdf}}
\pandocbounded{\includegraphics[keepaspectratio]{midterm-1_files/figure-latex/unnamed-chunk-3-49.pdf}}
\pandocbounded{\includegraphics[keepaspectratio]{midterm-1_files/figure-latex/unnamed-chunk-3-50.pdf}}
\pandocbounded{\includegraphics[keepaspectratio]{midterm-1_files/figure-latex/unnamed-chunk-3-51.pdf}}
\pandocbounded{\includegraphics[keepaspectratio]{midterm-1_files/figure-latex/unnamed-chunk-3-52.pdf}}
\pandocbounded{\includegraphics[keepaspectratio]{midterm-1_files/figure-latex/unnamed-chunk-3-53.pdf}}
\pandocbounded{\includegraphics[keepaspectratio]{midterm-1_files/figure-latex/unnamed-chunk-3-54.pdf}}
\pandocbounded{\includegraphics[keepaspectratio]{midterm-1_files/figure-latex/unnamed-chunk-3-55.pdf}}
\pandocbounded{\includegraphics[keepaspectratio]{midterm-1_files/figure-latex/unnamed-chunk-3-56.pdf}}
\pandocbounded{\includegraphics[keepaspectratio]{midterm-1_files/figure-latex/unnamed-chunk-3-57.pdf}}
\pandocbounded{\includegraphics[keepaspectratio]{midterm-1_files/figure-latex/unnamed-chunk-3-58.pdf}}
\pandocbounded{\includegraphics[keepaspectratio]{midterm-1_files/figure-latex/unnamed-chunk-3-59.pdf}}
\pandocbounded{\includegraphics[keepaspectratio]{midterm-1_files/figure-latex/unnamed-chunk-3-60.pdf}}
\pandocbounded{\includegraphics[keepaspectratio]{midterm-1_files/figure-latex/unnamed-chunk-3-61.pdf}}
\pandocbounded{\includegraphics[keepaspectratio]{midterm-1_files/figure-latex/unnamed-chunk-3-62.pdf}}
\pandocbounded{\includegraphics[keepaspectratio]{midterm-1_files/figure-latex/unnamed-chunk-3-63.pdf}}
\pandocbounded{\includegraphics[keepaspectratio]{midterm-1_files/figure-latex/unnamed-chunk-3-64.pdf}}
\pandocbounded{\includegraphics[keepaspectratio]{midterm-1_files/figure-latex/unnamed-chunk-3-65.pdf}}
\pandocbounded{\includegraphics[keepaspectratio]{midterm-1_files/figure-latex/unnamed-chunk-3-66.pdf}}
\pandocbounded{\includegraphics[keepaspectratio]{midterm-1_files/figure-latex/unnamed-chunk-3-67.pdf}}
\pandocbounded{\includegraphics[keepaspectratio]{midterm-1_files/figure-latex/unnamed-chunk-3-68.pdf}}
\pandocbounded{\includegraphics[keepaspectratio]{midterm-1_files/figure-latex/unnamed-chunk-3-69.pdf}}
\pandocbounded{\includegraphics[keepaspectratio]{midterm-1_files/figure-latex/unnamed-chunk-3-70.pdf}}
\pandocbounded{\includegraphics[keepaspectratio]{midterm-1_files/figure-latex/unnamed-chunk-3-71.pdf}}
\pandocbounded{\includegraphics[keepaspectratio]{midterm-1_files/figure-latex/unnamed-chunk-3-72.pdf}}
\pandocbounded{\includegraphics[keepaspectratio]{midterm-1_files/figure-latex/unnamed-chunk-3-73.pdf}}
\pandocbounded{\includegraphics[keepaspectratio]{midterm-1_files/figure-latex/unnamed-chunk-3-74.pdf}}
\pandocbounded{\includegraphics[keepaspectratio]{midterm-1_files/figure-latex/unnamed-chunk-3-75.pdf}}
\pandocbounded{\includegraphics[keepaspectratio]{midterm-1_files/figure-latex/unnamed-chunk-3-76.pdf}}
\pandocbounded{\includegraphics[keepaspectratio]{midterm-1_files/figure-latex/unnamed-chunk-3-77.pdf}}
\pandocbounded{\includegraphics[keepaspectratio]{midterm-1_files/figure-latex/unnamed-chunk-3-78.pdf}}
\pandocbounded{\includegraphics[keepaspectratio]{midterm-1_files/figure-latex/unnamed-chunk-3-79.pdf}}
\pandocbounded{\includegraphics[keepaspectratio]{midterm-1_files/figure-latex/unnamed-chunk-3-80.pdf}}
\pandocbounded{\includegraphics[keepaspectratio]{midterm-1_files/figure-latex/unnamed-chunk-3-81.pdf}}
\pandocbounded{\includegraphics[keepaspectratio]{midterm-1_files/figure-latex/unnamed-chunk-3-82.pdf}}
\pandocbounded{\includegraphics[keepaspectratio]{midterm-1_files/figure-latex/unnamed-chunk-3-83.pdf}}
\pandocbounded{\includegraphics[keepaspectratio]{midterm-1_files/figure-latex/unnamed-chunk-3-84.pdf}}
\pandocbounded{\includegraphics[keepaspectratio]{midterm-1_files/figure-latex/unnamed-chunk-3-85.pdf}}
\pandocbounded{\includegraphics[keepaspectratio]{midterm-1_files/figure-latex/unnamed-chunk-3-86.pdf}}
\pandocbounded{\includegraphics[keepaspectratio]{midterm-1_files/figure-latex/unnamed-chunk-3-87.pdf}}
\pandocbounded{\includegraphics[keepaspectratio]{midterm-1_files/figure-latex/unnamed-chunk-3-88.pdf}}
\pandocbounded{\includegraphics[keepaspectratio]{midterm-1_files/figure-latex/unnamed-chunk-3-89.pdf}}
\pandocbounded{\includegraphics[keepaspectratio]{midterm-1_files/figure-latex/unnamed-chunk-3-90.pdf}}
\pandocbounded{\includegraphics[keepaspectratio]{midterm-1_files/figure-latex/unnamed-chunk-3-91.pdf}}
\pandocbounded{\includegraphics[keepaspectratio]{midterm-1_files/figure-latex/unnamed-chunk-3-92.pdf}}
\pandocbounded{\includegraphics[keepaspectratio]{midterm-1_files/figure-latex/unnamed-chunk-3-93.pdf}}
\pandocbounded{\includegraphics[keepaspectratio]{midterm-1_files/figure-latex/unnamed-chunk-3-94.pdf}}
\pandocbounded{\includegraphics[keepaspectratio]{midterm-1_files/figure-latex/unnamed-chunk-3-95.pdf}}
\pandocbounded{\includegraphics[keepaspectratio]{midterm-1_files/figure-latex/unnamed-chunk-3-96.pdf}}
\pandocbounded{\includegraphics[keepaspectratio]{midterm-1_files/figure-latex/unnamed-chunk-3-97.pdf}}
\pandocbounded{\includegraphics[keepaspectratio]{midterm-1_files/figure-latex/unnamed-chunk-3-98.pdf}}
\pandocbounded{\includegraphics[keepaspectratio]{midterm-1_files/figure-latex/unnamed-chunk-3-99.pdf}}
\pandocbounded{\includegraphics[keepaspectratio]{midterm-1_files/figure-latex/unnamed-chunk-3-100.pdf}}

\begin{enumerate}
\def\labelenumi{\arabic{enumi}.}
\setcounter{enumi}{3}
\item
  (10 points) Above and Beyond. Choose either option below. You are
  welcome to explore both, but only one is required.

  \begin{itemize}
  \item
    Option 1: Explore the data a bit, and create a graphic that uses (or
    incorporates) different information than what was used above. Some
    notes:

    \begin{itemize}
    \tightlist
    \item
      Create an additional graphic that incorporates at least one
      additional variable not previously used (this should add to the
      graphic in part 1). The additional data should be drawn from a
      different dataset (function call) than the original graphic used.
      These two (or more) datasets may need to be joined appropriately.
    \item
      You can either add more information to the plot above, or create a
      different plot.
    \item
      Formatting, labelling, etc. are all important here too.
    \item
      Adding marginal densities or other ``bells and whistles'' might
      offer additional insight.
    \item
      This graphic should be included at the end of the report (after
      the more detailed explanations).
    \item
      You should include a brief description of the graphic
      (highlighting the different/additional information used).
    \end{itemize}
  \item
    Option 2: If the NBA were to incorporate a 4-point shot, where would
    you draw a 4-point arc? Some notes:

    \begin{itemize}
    \tightlist
    \item
      You likely should base your decision at least partly on proportion
      of shots made from a given range. You might consider an expected
      value calculation here.
    \item
      Your arc can be shaped as you see fit; simple arcs are sufficient
      for this exploration.
    \item
      Provide an example of a consequence (positive or negative) if a
      4-point shot was incorporated. (e.g., ``my\_favorite\_player's
      season point total would increase by x\%'').
    \item
      You do not need to create a plot representing your arc, though you
      are welcome to do so!
    \end{itemize}
  \end{itemize}
\end{enumerate}

\subsubsection{Analysis/Reasoning:}\label{analysisreasoning}

The NBA three point line measures at 23.75 feet straight out from the
base line. According to
\url{https://www.statmuse.com/nba/ask/pelicans-3-point-shooting-percentage},
the Pelican's shooting percentage for three pointers in 2025-2026 is
34.7\%. Using that percentage, an expected value for threes would be 3 *
34.7\% = 1.04. Based on that, the average make percentage for an
imagined 4 point shot would have to be 1.04 / 4 = 26.03\%. Then, to
determine where to put the new 4 point line, I calculated how far out
the Pelicans were making roughly 26.03\% of their shots.

\begin{Shaded}
\begin{Highlighting}[]
\NormalTok{part4data }\OtherTok{\textless{}{-}}\NormalTok{ baby\_data }\SpecialCharTok{|\textgreater{}} \FunctionTok{filter}\NormalTok{(coordinate\_y\_raw }\SpecialCharTok{\textgreater{}=} \DecValTok{28}\NormalTok{)}
\NormalTok{makeCount }\OtherTok{\textless{}{-}} \FunctionTok{sum}\NormalTok{(part4data}\SpecialCharTok{$}\NormalTok{make\_miss }\SpecialCharTok{==} \StringTok{"Make"}\NormalTok{)}
\NormalTok{percent28 }\OtherTok{\textless{}{-}}\NormalTok{ makeCount }\SpecialCharTok{/} \FunctionTok{nrow}\NormalTok{(part4data)}
\NormalTok{percent28}
\end{Highlighting}
\end{Shaded}

\begin{verbatim}
## [1] 0.2295633
\end{verbatim}

\begin{Shaded}
\begin{Highlighting}[]
\NormalTok{part4data }\OtherTok{\textless{}{-}}\NormalTok{ baby\_data }\SpecialCharTok{|\textgreater{}} \FunctionTok{filter}\NormalTok{(coordinate\_y\_raw }\SpecialCharTok{\textgreater{}=} \DecValTok{27}\NormalTok{)}
\NormalTok{makeCount }\OtherTok{\textless{}{-}} \FunctionTok{sum}\NormalTok{(part4data}\SpecialCharTok{$}\NormalTok{make\_miss }\SpecialCharTok{==} \StringTok{"Make"}\NormalTok{)}
\NormalTok{percent27 }\OtherTok{\textless{}{-}}\NormalTok{ makeCount }\SpecialCharTok{/} \FunctionTok{nrow}\NormalTok{(part4data)}
\NormalTok{percent27}
\end{Highlighting}
\end{Shaded}

\begin{verbatim}
## [1] 0.2859477
\end{verbatim}

Based on the above code, I've determined the distance for a 4 point line
should be 27.5 feet out from the base line. The Pelicans make roughly
28.59\% of their shots at or beyond 27 feet from the base line, but only
22.96\% of their shots at or beyond 28 feet. This makes 27.5 feet a
happy medium that should get the make percentage close to 26\%. Doing
some simple math, 27.5 - 23.75 = 3.75, the 4 point line would follow the
current 3 point line, but be 3.75 feet further from the basket.

\subsubsection{Consequences:}\label{consequences}

The addition of a 4 point line at 27.5 feet from the base line would
have positive results for teams with players who consistently score deep
threes. For anyone who can consistently score beyond 27.5 feet, all
those former three point shots become four point shots, increasing their
overall points per game which in turn increases their team's overall
scoring output, causing their team to be more successful in games.

Assuming my favorite player on the New Orleans Pelicans is Trey Murphy
III, who currently has a 38.3 three point percentage, and, according to
\url{https://thebirdwrites.com/triggad-trey-murphy-iiis-elite-shooting-sparks-pelicans/},
41.8\% of his threes are from 25-29 feet. With that kind of range, he
would definitely be making a significant amount of 4 pointers,
increasing his average points per game and giving his team a huge
offensive advantage.

\subsection{The Deliverables}\label{the-deliverables}

\begin{enumerate}
\def\labelenumi{\arabic{enumi}.}
\tightlist
\item
  Upload your report and code file(s) to GitHub by 5:00pm on Friday,
  April 3.
\item
  Submit (on Canvas) your report, code, and link to your GitHub
  repository by 5:00pm on Friday, April 3.
\end{enumerate}

\end{document}
